\chapter{René Frédéric Thom --- Fields Medal 1958}\label{ch1958renethom:renuxe9-fruxe9duxe9ric-thom-fields-medal-1958-a-comprehensive-report}

\begin{center}\rule{0.5\linewidth}{0.5pt}\end{center}

\textbf{Prepared:} July 2026 \textbf{Subject:} René Frédéric Thom (2
September 1923 -- 25 October 2002) \textbf{Position:} Professor,
Institut des Hautes Études Scientifiques (IHÉS) \textbf{Doctoral
Advisor:} Henri Cartan (Université de Paris / École Normale Supérieure)
\textbf{Fields Medal Citation:} \emph{``The creator of cobordism theory,
a central concept in algebraic topology that classifies manifolds up to
the existence of a common boundary. His work on the Thom transversality
theorem, Thom spaces, and characteristic classes revolutionized
differential topology and laid the foundation for modern singularity
theory and catastrophe theory.''}

\begin{center}\rule{0.5\linewidth}{0.5pt}\end{center}

\section{1. Introduction}\label{ch1958renethom:introduction}

René Thom was awarded the Fields Medal in 1958 at the International
Congress of Mathematicians in Edinburgh for his substantial
contributions to algebraic and differential topology, particularly the
creation of \textbf{cobordism theory}, the \textbf{Thom transversality
theorem}, and the theory of \textbf{Thom spaces} and
\textbf{characteristic classes}. His work transformed the study of
manifolds and singularities.

\begin{prereqs} Differential topology (manifolds, transversality), basic singularity theory, multivariable calculus. \end{prereqs}

\begin{objectives} State the classification theorem for smooth manifolds of dimension $< 5$; explain the cobordism theory; describe the seven elementary catastrophes. \end{objectives}


Thom's major contributions include:

\begin{itemize}
\tightlist
\item
 \textbf{Cobordism theory:} A classification of manifolds up to the
 relation of being the boundary of a higher-dimensional manifold. Thom
 completely computed the cobordism rings and showed that the cobordism
 class of a manifold is determined by its characteristic numbers.
\item
 \textbf{The Thom transversality theorem:} A fundamental result stating
 that any continuous map can be approximated by one that is transverse
 to a given submanifold, a cornerstone of differential topology.
\item
 \textbf{The Thom isomorphism theorem:} A duality between the
 cohomology of a vector bundle and that of its base space, giving rise
 to the \textbf{Thom class} and \textbf{Euler class}.
\item
 \textbf{Characteristic classes:} A new understanding of
 Stiefel--Whitney classes using Steenrod operations applied to the Thom
 class.
\item
 \textbf{Singularity theory:} The classification of singularities of
 differentiable maps, leading to his well-known \textbf{catastrophe
 theory}.
\item
 \textbf{Stratified sets and the Thom--Mather isotopy theorem:} The
 theory of Whitney stratified sets and the topological stability of
 maps.
\end{itemize}

This report provides a comprehensive overview of Thom's life and his
monumental contributions to mathematics.

\begin{center}\rule{0.5\linewidth}{0.5pt}\end{center}

\section{2. Biographical Background}\label{ch1958renethom:biographical-background}

\subsection{2.1 Early Life and
Education}\label{ch1958renethom:early-life-and-education}

René Frédéric Thom was born on 2 September 1923 in Montbéliard, Doubs,
France. He showed academic potential early, winning a scholarship to
Lycée Saint-Louis in Paris. He entered the \textbf{École Normale
Supérieure (ENS)} in 1943, where he studied under Henri Cartan.

\subsection{2.2 Academic Career}\label{ch1958renethom:academic-career}

\begin{itemize}
\tightlist
\item
 \textbf{1946:} Graduated from ENS.
\item
 \textbf{1946--1951:} CNRS researcher at the University of Strasbourg,
 working with Henri Cartan, Charles Ehresmann, and Jean-Louis Koszul.
\item
 \textbf{1951:} Doctorate from the University of Paris, supervised by
 Henri Cartan. Thesis: \emph{``Espaces fibrés en sphères et carrés de
 Steenrod''} (Fiber spaces in spheres and Steenrod squares).
\item
 \textbf{1951--1952:} Rockefeller Fellow in the United States, visiting
 Princeton (where he met Einstein, Weyl, and Steenrod) and Stanford.
\item
 \textbf{1953--1954:} Professor at the University of Grenoble.
\item
 \textbf{1954--1963:} Professor at the University of Strasbourg.
\item
 \textbf{1957:} Became full professor at age 34.
\item
 \textbf{1958:} Fields Medal, ICM Edinburgh.
\item
 \textbf{1963--1990:} Professor at the \textbf{Institut des Hautes
 Études Scientifiques (IHÉS)}, Bures-sur-Yvette.
\item
 \textbf{1970:} Brouwer Medal.
\item
 \textbf{1990:} Emeritus Professor, IHÉS.
\item
 \textbf{2002:} Died at age 79 in Bures-sur-Yvette.
\end{itemize}

\subsection{2.3 Doctoral Students}\label{ch1958renethom:doctoral-students}

Thom's doctoral students include \textbf{David Trotman} and several
other mathematicians who developed singularity theory and stratified
spaces.

\begin{center}\rule{0.5\linewidth}{0.5pt}\end{center}

\section{3. The Thom Space and Thom
Isomorphism}\label{ch1958renethom:the-thom-space-and-thom-isomorphism}

\subsection{3.1 Definition}\label{ch1958renethom:definition}

Let \(\xi: E \to B\) be a real vector bundle of rank \(n\). The
\textbf{Thom space} \(T(\xi)\) is defined as the one-point
compactification of the total space \(E\), or equivalently as
\(D(E)/S(E)\), where \(D(E)\) is the disk bundle and \(S(E)\) the sphere
bundle.

\subsection{3.2 The Thom Isomorphism}\label{ch1958renethom:the-thom-isomorphism}

The \textbf{Thom isomorphism theorem} states that for an oriented vector
bundle \(\xi\) of rank \(n\):

\[H^q(B) \cong H^{q+n}(D(E), S(E); \mathbb{Z})\]

or equivalently, there is an isomorphism:

\[\phi: H^q(B) \xrightarrow{\cong} \widetilde{H}^{q+n}(T(\xi))\]

given by \(\phi(x) = p^*x \cup u\), where \(u \in H^n(D(E), S(E))\) is
the \textbf{Thom class}.

\subsection{3.3 The Thom Class}\label{ch1958renethom:the-thom-class}

The \textbf{Thom class} \(u \in H^n(D(E), S(E))\) is the unique
cohomology class that restricts to the orientation class on each fiber.
For the tangent bundle of an oriented manifold \(M^n\), the Thom class
gives the \textbf{Euler class} \(e(M) = s^*u\) by pulling back via the
zero section.

\subsection{3.4 The Euler Class}\label{ch1958renethom:the-euler-class}

The \textbf{Euler class} \(e(\xi) \in H^n(B)\) is the pullback of the
Thom class via the zero section. For the tangent bundle of an oriented
\(n\)-manifold \(M\), the Euler class satisfies \(e(M)[M] = \chi(M)\),
where \(\chi(M)\) is the Euler characteristic.

\subsection{3.5 The Thom Spectrum}\label{ch1958renethom:the-thom-spectrum}

The \textbf{Thom spectrum \(MO\)} is the spectrum whose \(n\)th space is
the Thom space \(T(\gamma^n)\) of the universal vector bundle over
\(BO(n)\):

\[MO_n = T(\gamma^n)\]

The structure maps \(\Sigma MO_n \to MO_{n+1}\) come from the natural
inclusion \(BO(n) \hookrightarrow BO(n+1)\). The homotopy groups of
\(MO\) are the unoriented cobordism groups:

\[\pi_n(MO) \cong \Omega_n^{\text{O}}\]

This connection between topology and manifold theory is the cornerstone
of cobordism theory.

\subsection{3.6 Complex Thom Spectra}\label{ch1958renethom:complex-thom-spectra}

Similarly, the \textbf{complex Thom spectrum \(MU\)} has \(n\)th space
\(T(\gamma^n_{\mathbb{C}})\) where \(\gamma^n_{\mathbb{C}}\) is the
universal complex \(n\)-plane bundle over \(BU(n)\). The homotopy groups
\(\pi_*(MU)\) give the \textbf{complex cobordism ring}, which was
completely computed by \textbf{Milnor} and \textbf{Novikov}:

\[\pi_*(MU) \cong \mathbb{Z}[x_2, x_4, x_6, \ldots]\]

The generators \(x_{2n}\) correspond to the cobordism classes of
\(\mathbb{CP}^n\).

\begin{center}\rule{0.5\linewidth}{0.5pt}\end{center}

\section{4. Cobordism Theory}\label{ch1958renethom:cobordism-theory}

\subsection{4.1 Definition}\label{ch1958renethom:definition-1}

Two closed \(n\)-manifolds \(M_1\) and \(M_2\) are \textbf{cobordant} if
their disjoint union is the boundary of a compact \((n+1)\)-manifold
\(W\). The \textbf{bordism group} \(\Omega_n\) consists of
\(n\)-manifolds modulo the cobordism relation.

\subsection{4.2 The Pontryagin--Thom
Construction}\label{ch1958renethom:the-pontryaginthom-construction}

The \textbf{Pontryagin--Thom construction} is the central technical tool
of cobordism theory. A closed \(n\)-manifold
\(M \subset \mathbb{R}^{n+k}\) has a normal bundle classified by a map
\(\nu: M \to BO(k)\). The \textbf{Thom construction} produces a map:

\[S^{n+k} \to T(\gamma^k)\]

where \(\gamma^k\) is the universal bundle over \(BO(k)\) and
\(T(\gamma^k)\) is the Thom space. This gives a bijection:

\[\Omega_n^{\text{SO}} \cong \pi_{n+k}(MSO(k)) \cong \lim_{k\to\infty} \pi_{n+k}(MSO(k))\]

where \(MSO(k)\) is the Thom space of the universal oriented bundle.

\subsection{4.3 Thom's Computation}\label{ch1958renethom:thoms-computation}

Thom completely computed the \textbf{unoriented cobordism ring}
\(\Omega_*^{\text{O}}\):

\[\Omega_*^{\text{O}} \cong \mathbb{F}_2[x_2, x_4, x_5, x_6, x_8, \ldots]\]

where the generators \(x_i\) correspond to the cobordism classes of real
projective spaces. The ring is a polynomial ring over \(\mathbb{F}_2\)
with one generator in each degree \(i\) not of the form \(2^k - 1\).

\subsection{4.4 The Oriented Cobordism
Ring}\label{ch1958renethom:the-oriented-cobordism-ring}

For oriented cobordism, Thom computed:

\[\Omega_*^{\text{SO}} \otimes \mathbb{Q} \cong \mathbb{Q}[x_4, x_8, x_{12}, \ldots]\]

The structure of \(\Omega_*^{\text{SO}}\) over \(\mathbb{Z}\) was fully
determined later by \textbf{Wall} and \textbf{Milnor}, requiring the
computation of torsion.

\subsection{4.5 Characteristic Numbers}\label{ch1958renethom:characteristic-numbers}

Two manifolds are cobordant if and only if all their
\textbf{Stiefel--Whitney numbers} (for unoriented cobordism) or
\textbf{Pontryagin numbers} (for oriented cobordism) coincide. This is
Thom's well-known theorem:

\begin{quote}
\textbf{Thom's Cobordism Theorem:} \(M\) is a boundary if and only if
all Stiefel--Whitney numbers of \(M\) vanish.
\end{quote}

\subsection{4.6 Examples of Cobordant
Manifolds}\label{ch1958renethom:examples-of-cobordant-manifolds}

\begin{itemize}
\tightlist
\item
 \(\mathbb{RP}^{2n}\) has Stiefel--Whitney number
 \(w_{2n}[\mathbb{RP}^{2n}] \neq 0\), so it is not a boundary.
\item
 \(\mathbb{RP}^{2n+1}\) has all Stiefel--Whitney numbers zero, so it
 bounds a \((2n+2)\)-manifold.
\item
 \(\mathbb{CP}^n\) has all Stiefel--Whitney numbers zero, but its
 Pontryagin numbers may be non-zero.
\item
 The \textbf{Klein bottle} is cobordant to the torus (they have the
 same Stiefel--Whitney numbers).
\end{itemize}

\subsection{4.7 The Thom Polynomials}\label{ch1958renethom:the-thom-polynomials}

The relation between cobordism and Stiefel--Whitney numbers is given by
\textbf{Thom polynomials}. For a given \(\omega\) (a partition), the
Stiefel--Whitney number \(w_\omega[M]\) is a homomorphism from
\(\Omega_n^{\text{O}}\) to \(\mathbb{F}_2\). The \textbf{duality} is
that these homomorphisms separate cobordism classes.

\subsection{4.8 Oriented Cobordism and Pontryagin
Numbers}\label{ch1958renethom:oriented-cobordism-and-pontryagin-numbers}

For oriented cobordism, the invariants are \textbf{Pontryagin numbers}:

\[p_\omega[M] = \langle p_{i_1} \cup \cdots \cup p_{i_k}, [M]\rangle\]

where \(p_i\) are the Pontryagin classes of \(TM\). Two oriented
manifolds are oriented-cobordant if and only if all their Pontryagin
numbers coincide. Over \(\mathbb{Q}\), the oriented cobordism ring is:

\[\Omega_*^{\text{SO}} \otimes \mathbb{Q} \cong \mathbb{Q}[x_4, x_8, x_{12}, \ldots]\]

where \(x_{4k} = [\mathbb{CP}^{2k}]\).

\begin{center}\rule{0.5\linewidth}{0.5pt}\end{center}

\section{5. The Thom Transversality
Theorem}\label{ch1958renethom:the-thom-transversality-theorem}

\subsection{5.1 Statement}\label{ch1958renethom:statement}

The \textbf{Thom transversality theorem} is a foundational result in
differential topology:

\begin{quote}
\textbf{Thom Transversality Theorem:} Let \(f: M \to N\) be a smooth map
of manifolds and let \(X \subset N\) be a submanifold. Then \(f\) can be
approximated arbitrarily closely by a smooth map \(g\) that is
\textbf{transverse} to \(X\).
\end{quote}

\subsection{5.2 Transversality}\label{ch1958renethom:transversality}

A map \(f: M \to N\) is \textbf{transverse} to a submanifold
\(X \subset N\) at \(x \in M\) if either \(f(x) \notin X\) or:

\[T_{f(x)}N = df(T_xM) + T_{f(x)}X\]

If \(f\) is transverse to \(X\), then \(f^{-1}(X) \subset M\) is a
submanifold of codimension equal to the codimension of \(X\) in \(N\).

\subsection{5.3 Applications}\label{ch1958renethom:applications}

The transversality theorem is essential for: - Defining
\textbf{intersection theory} on manifolds. - Constructing
\textbf{Thom--Pontryagin constructions}. - Defining \textbf{Euler
classes} and \textbf{Chern classes} via generic sections. - The theory
of \textbf{moduli spaces} in gauge theory. - \textbf{Morse theory} and
the \textbf{h-principle}.

\subsection{5.4 The Multijet Transversality
Theorem}\label{ch1958renethom:the-multijet-transversality-theorem}

Thom also proved the \textbf{multijet transversality theorem}, which
guarantees that for a generic map, the \(k\)-fold jet extension is
transverse to a given submanifold of the jet bundle. This is the key to
the classification of singularities.

\begin{center}\rule{0.5\linewidth}{0.5pt}\end{center}

\section{6. Characteristic Classes and Steenrod
Operations}\label{ch1958renethom:characteristic-classes-and-steenrod-operations}

\subsection{6.1 Stiefel--Whitney Classes via the Thom
Class}\label{ch1958renethom:stiefelwhitney-classes-via-the-thom-class}

Thom discovered a fundamental relationship between \textbf{Steenrod
operations} and \textbf{Stiefel--Whitney classes}. For a vector bundle
\(\xi\) with Thom class \(u \in H^n(T(\xi); \mathbb{F}_2)\):

\[w_i(\xi) = \Phi^{-1}(Sq^i(u))\]

where \(\Phi\) is the Thom isomorphism and \(Sq^i\) are the Steenrod
squares. This gives a direct cohomological definition of
Stiefel--Whitney classes.

\subsection{6.2 The Wu Formula}\label{ch1958renethom:the-wu-formula}

The \textbf{Wu formula} expresses the Stiefel--Whitney classes of a
manifold in terms of the Steenrod operations on its cohomology:

\[w_i(M) = \sum_{j=0}^i Sq^j(v_{i-j})\]

where \(v_k \in H^k(M; \mathbb{F}_2)\) are the \textbf{Wu classes}
defined by:

\[\langle Sq^k(x), [M]\rangle = \langle v_k \cup x, [M]\rangle\]

This formula, derived using Thom's insight about Steenrod squares and
the Thom class, is a cornerstone of characteristic class theory.

\subsection{6.3 The Thom--Porteous
Formula}\label{ch1958renethom:the-thomporteous-formula}

The \textbf{Thom--Porteous formula} (generalized by Porteous) expresses
the fundamental class of a degeneracy locus in terms of Chern classes.
For a map \(\sigma: E \to F\) of vector bundles, the locus where
\(\text{rank}(\sigma) \leq r\) has cohomology class given by a
determinant in Chern classes.

\begin{center}\rule{0.5\linewidth}{0.5pt}\end{center}

\section{7. Singularity Theory}\label{ch1958renethom:singularity-theory}

\subsection{7.1 Germs of Functions}\label{ch1958renethom:germs-of-functions}

Thom initiated the systematic study of \textbf{singularities of
differentiable maps}. A \textbf{germ} of a smooth function
\(f: \mathbb{R}^n \to \mathbb{R}\) at \(0\) has a singularity if
\(df(0) = 0\). The type of singularity is classified by the
\textbf{Taylor expansion} up to smooth coordinate changes.

\subsection{7.2 The Classification of Simple
Singularities}\label{ch1958renethom:the-classification-of-simple-singularities}

Thom classified the \textbf{simple singularities} (also called
\textbf{ADE singularities}): - \textbf{\(A_k\)}: \(x^{k+1} + y^2 + z^2\)
(fold, cusp, etc.) - \textbf{\(D_k\)}: \(x^{k-1} + xy^2 + z^2\) -
\textbf{\(E_6\)}: \(x^3 + y^4 + z^2\) - \textbf{\(E_7\)}:
\(x^3 + xy^3 + z^2\) - \textbf{\(E_8\)}: \(x^3 + y^5 + z^2\)

This classification is intimately connected to Dynkin diagrams and the
theory of Lie groups.

\subsection{7.3 The Thom-Boardman
Singularities}\label{ch1958renethom:the-thom-boardman-singularities}

The \textbf{Thom--Boardman stratification} classifies singularities of
maps \(f: M \to N\) by the ranks of their derivatives. The
Thom--Boardman symbols \(\Sigma^{i,j,k,\ldots}\) encode the pattern of
rank drops in successive derivatives.

\subsection{7.4 Finite Determinacy}\label{ch1958renethom:finite-determinacy}

A germ \(f\) is \textbf{\(k\)-determined} if any other germ with the
same \(k\)-jet is right-equivalent to \$f). Thom's work on finite
determinacy established that: - Isolated singularities are finitely
determined. - The determinacy degree can be computed from the
\textbf{Milnor number} and the \textbf{Tjurina number}.

\begin{center}\rule{0.5\linewidth}{0.5pt}\end{center}

\section{8. Catastrophe Theory}\label{ch1958renethom:catastrophe-theory}

\subsection{8.1 Motivation}\label{ch1958renethom:motivation}

\textbf{Catastrophe theory} (developed by Thom in the 1960s, published
in \emph{Structural Stability and Morphogenesis}, 1972) studies how
continuous changes in parameters can produce discontinuous changes in
behavior --- so-called \textbf{catastrophes}.

\subsection{8.2 The Seven Elementary
Catastrophes}\label{ch1958renethom:the-seven-elementary-catastrophes}

Thom classified seven elementary catastrophes for systems with up to
four control parameters and two state variables:

{\def\LTcaptype{none} % do not increment counter
\begin{longtable}[]{@{}lll@{}}
\toprule\noalign{}
Name & Codimension & Universal unfolding \\
\midrule\noalign{}
\endhead
\bottomrule\noalign{}
\endlastfoot
\textbf{Fold} & 1 & \(x^3 + ux\) \\
\textbf{Cusp} & 2 & \(x^4 + ux^2 + vx\) \\
\textbf{Swallowtail} & 3 & \(x^5 + ux^3 + vx^2 + wx\) \\
\textbf{Butterfly} & 4 & \(x^6 + ux^4 + vx^3 + wx^2 + zx\) \\
\textbf{Hyperbolic umbilic} & 3 & \(x^3 + y^3 + uxy + vx + wy\) \\
\textbf{Elliptic umbilic} & 3 & \(x^3 - xy^2 + u(x^2+y^2) + vx + wy\) \\
\textbf{Parabolic umbilic} & 4 &
\(x^2y + y^4 + ux^2 + vy^2 + wx + zy\) \\
\end{longtable}
}

\subsection{8.3 The Classification
Theorem}\label{ch1958renethom:the-classification-theorem}

Thom's classification theorem states: for a generic family of smooth
functions depending on at most four parameters, only the seven
elementary catastrophes appear as stable singularities.

\subsection{8.4 Applications}\label{ch1958renethom:applications-1}

Catastrophe theory has been applied to: - \textbf{Biology:} Population
dynamics, morphogenesis. - \textbf{Physics:} Phase transitions, optics
(caustics). - \textbf{Engineering:} Structural stability. -
\textbf{Economics:} Market crashes and discontinuities. -
\textbf{Psychology:} Models of behavioral change (Zeeman's
applications).

\subsection{8.5 Controversy and Legacy}\label{ch1958renethom:controversy-and-legacy}

Catastrophe theory generated great popular interest but also
controversy. Critics questioned the validity of some applications,
particularly in the social sciences. However, the underlying mathematics
--- singularity theory and the classification of stable maps --- remains
a central part of differential topology.

\begin{center}\rule{0.5\linewidth}{0.5pt}\end{center}

\section{9. Stratified Sets}\label{ch1958renethom:stratified-sets}

\subsection{9.1 Whitney Stratifications}\label{ch1958renethom:whitney-stratifications}

A \textbf{Whitney stratification} of a space \(X\) is a decomposition
into smooth manifolds (strata) satisfying Whitney's conditions (a) and
(b). These conditions control the behavior of the stratification near
the boundary.

\subsection{9.2 The Thom--Mather Isotopy
Theorem}\label{ch1958renethom:the-thommather-isotopy-theorem}

The \textbf{Thom--Mather isotopy theorem} is a fundamental result in the
theory of stratified spaces:

\begin{quote}
\textbf{Thom--Mather Theorem:} A proper submersion between Whitney
stratified sets that is \textbf{controlled} (satisfies the Thom \(A_f\)
condition or Mather's control conditions) is a topological fiber bundle.
\end{quote}

\subsection{9.3 The First Isotopy Lemma}\label{ch1958renethom:the-first-isotopy-lemma}

The \textbf{first isotopy lemma} states that a proper stratified
submersion is locally trivial: the topological type of the fibers does
not change along a stratum. This is the key tool for understanding the
topology of singular spaces.

\subsection{9.4 The Second Isotopy
Lemma}\label{ch1958renethom:the-second-isotopy-lemma}

The \textbf{second isotopy lemma} (Mather) states that the tube system
of a Whitney stratification can be radially deformed, giving a local
product structure near each stratum.

\begin{center}\rule{0.5\linewidth}{0.5pt}\end{center}

\section{10. Dold--Thom Theorem}\label{ch1958renethom:doldthom-theorem}

\subsection{10.1 Statement}\label{ch1958renethom:statement-1}

The \textbf{Dold--Thom theorem} (proved with Albrecht Dold) relates the
homotopy groups of the symmetric product \(SP(X)\) to the homology
groups of \(X\):

\[\pi_n(SP(X)) \cong H_n(X; \mathbb{Z})\]

where \(SP(X)\) is the infinite symmetric product of \(X\).

\subsection{10.2 Significance}\label{ch1958renethom:significance}

This theorem provides a purely homotopy-theoretic definition of homology
groups, showing that homology arises naturally from the abelianization
of the fundamental group of the symmetric product.

\subsection{10.3 Applications}\label{ch1958renethom:applications-2}

The Dold--Thom theorem has been used in: - \textbf{Algebraic K-theory}
(the Quillen Q-construction). - \textbf{Motivic homotopy theory} (the
Morel--Voevodsky \(\mathbb{A}^1\)-homotopy category). - \textbf{The
theory of configuration spaces.}

\begin{center}\rule{0.5\linewidth}{0.5pt}\end{center}

\section{11. The Thom Conjecture}\label{ch1958renethom:the-thom-conjecture}

\subsection{11.1 Statement}\label{ch1958renethom:statement-2}

The \textbf{Thom conjecture} (proved by Kronheimer and Mrowka using
Seiberg--Witten theory) states:

\begin{quote}
\textbf{Thom Conjecture (now theorem):} The complex curves of minimal
genus in a complex projective surface are algebraic. More precisely, for
a given homology class in \(\mathbb{CP}^2\), the minimal genus of an
embedded surface representing that class is realized by an algebraic
curve.
\end{quote}

\subsection{11.2 Generalizations}\label{ch1958renethom:generalizations}

The Thom conjecture has been generalized to: - \textbf{Symplectic
manifolds} (the symplectic Thom conjecture, proved by Ozsváth and
Szabó). - \textbf{Kähler surfaces} using the adjunction inequality. -
\textbf{4-manifolds} via Seiberg--Witten theory.

\begin{center}\rule{0.5\linewidth}{0.5pt}\end{center}

\section{12. The Porteous--Thom Formula}\label{ch1958renethom:the-porteousthom-formula}

\subsection{12.1 Statement}\label{ch1958renethom:statement-3}

The \textbf{Porteous--Thom formula} (developed with Ian Porteous) gives
the cohomology class of the degeneracy locus of a map of vector bundles.
For a bundle map \(\sigma: E \to F\) of ranks \(n\) and \(m\) with
\(n \leq m\):

\[D_r(\sigma) = \{\,x \in X : \text{rank}(\sigma_x) \leq r\,\}\]

has fundamental class:

\[[D_r(\sigma)] = \Delta_{m-r}(c(E^* \otimes F))\]

where \(\Delta_k\) denotes the \(k \times k\) determinant in Chern class
expressions.

\subsection{12.2 Applications}\label{ch1958renethom:applications-3}

This formula is essential in: - \textbf{Intersection theory} on
Grassmannians. - \textbf{The theory of Schubert varieties}. -
\textbf{Enumerative geometry}. - \textbf{The computation of
characteristic numbers} of vector bundle maps.

\begin{center}\rule{0.5\linewidth}{0.5pt}\end{center}

\section{13. Thom's Work on the
h-Principle}\label{ch1958renethom:thoms-work-on-the-h-principle}

\subsection{13.1 The h-Principle}\label{ch1958renethom:the-h-principle}

The \textbf{h-principle} (homotopy principle, developed by Gromov)
states that the existence of a solution to a partial differential
relation can be reduced to a purely topological condition. Thom's
transversality theorem and the theory of jets provided the foundations
for the h-principle.

\subsection{13.2 Thom's Formulation}\label{ch1958renethom:thoms-formulation}

Thom showed that the \textbf{jet extension} of a smooth map is
transverse to a given submanifold of the jet bundle for a generic map.
This implies that solutions to certain differential relations exist if
and only if the corresponding topology of the jet bundle allows it.

\subsection{13.3 Applications}\label{ch1958renethom:applications-4}

The h-principle has been applied to: - The \textbf{Nash embedding
theorem} (M. Gromov's proof). - \textbf{Symplectic geometry} (the
existence of Lagrangian immersions). - \textbf{Contact geometry}
(Legendrian knots). - \textbf{Immersions and embeddings} of manifolds.

\begin{center}\rule{0.5\linewidth}{0.5pt}\end{center}

\section{14. The Thom--Milnor Form}\label{ch1958renethom:the-thommilnor-form}

\subsection{14.1 Definition}\label{ch1958renethom:definition-2}

The \textbf{Thom--Milnor form} (or \textbf{Thom form}) is a differential
form representing the Thom class of an oriented vector bundle. It is a
closed \(n\)-form on the total space with compact vertical support,
integrating to \(1\) on each fiber.

\subsection{14.2 Construction}\label{ch1958renethom:construction}

Given a connection on an oriented vector bundle \(\xi\) with curvature
\(\Omega\), the Thom form can be constructed using the
\textbf{Mathai--Quillen formalism}:

\[\tau(\xi) = \frac{1}{(2\pi)^{n/2}} \exp\left(-\frac{|v|^2}{2}\right) \wedge \sum_{k \geq 0} \frac{\text{Pf}(\Omega)}{(2\pi)^k}\]

where \(v\) is the tautological section of the bundle and Pf is the
Pfaffian.

\subsection{14.3 Applications}\label{ch1958renethom:applications-5}

The Thom form is essential in: - \textbf{Chern--Weil theory} and
characteristic classes. - \textbf{Index theory} (the Atiyah--Singer
theorem). - \textbf{Localization formulas} in equivariant cohomology. -
\textbf{Axiomatic quantum field theory.}

\begin{center}\rule{0.5\linewidth}{0.5pt}\end{center}

\section{15. Thom's Work on Foliations}\label{ch1958renethom:thoms-work-on-foliations}

\subsection{15.1 The Thom--Haefliger
Theory}\label{ch1958renethom:the-thomhaefliger-theory}

Thom contributed to the theory of \textbf{foliations}, working with
André Haefliger. A foliation of codimension \(q\) on a manifold \(M\) is
a decomposition into immersed submanifolds of codimension \(q\).

\subsection{15.2 Characteristic Classes of
Foliations}\label{ch1958renethom:characteristic-classes-of-foliations}

Thom's work on the classification of foliations led to the theory of
\textbf{Haefliger structures} and \textbf{simplicial classifying spaces}
\(B\Gamma_q\). The \textbf{Godbillon--Vey} class is a secondary
characteristic class of foliations.

\subsection{15.3 The Cobordism of
Foliations}\label{ch1958renethom:the-cobordism-of-foliations}

Thom's cobordism theory was extended to foliated manifolds, giving rise
to \textbf{foliated cobordism} groups and the
\textbf{Haefliger--Thurston} theorem on the classification of
foliations.

\begin{center}\rule{0.5\linewidth}{0.5pt}\end{center}

\section{16. Thom's Work on Dynamical
Systems}\label{ch1958renethom:thoms-work-on-dynamical-systems}

\subsection{16.1 Structural Stability}\label{ch1958renethom:structural-stability}

Thom's work on structural stability of dynamical systems built on the
earlier work of Andronov and Pontryagin. He studied how the qualitative
behavior of a dynamical system changes under small perturbations.

\subsection{16.2 The Andronov--Hopf--Thom
Theorem}\label{ch1958renethom:the-andronovhopfthom-theorem}

The \textbf{Andronov--Hopf bifurcation} theorem describes the birth of
limit cycles from equilibria as a parameter is varied. Thom's
contribution was to understand this and other bifurcations in the
framework of singularity theory.

\subsection{16.3 Generic Properties}\label{ch1958renethom:generic-properties}

Thom proved that certain properties of vector fields are
\textbf{generic} (hold for an open dense set): - Hyperbolic equilibria
and periodic orbits. - Transverse intersections of stable and unstable
manifolds. - The Kupka--Smale property.

\begin{center}\rule{0.5\linewidth}{0.5pt}\end{center}

\section{17. The Thom--Sullivan Theory}\label{ch1958renethom:the-thomsullivan-theory}

\subsection{17.1 Rational Homotopy
Theory}\label{ch1958renethom:rational-homotopy-theory}

The \textbf{Thom--Sullivan theory} (developed with Dennis Sullivan)
connects rational homotopy theory to differential forms. The key idea is
that the \textbf{de Rham complex} of a manifold encodes information
about its rational homotopy type.

\subsection{17.2 Minimal Models}\label{ch1958renethom:minimal-models}

The \textbf{Sullivan minimal model} of a space is a free commutative
differential graded algebra that encodes the rational homotopy groups.
Thom's work on forms and cohomology provided the geometric foundation
for this algebraic construction.

\begin{center}\rule{0.5\linewidth}{0.5pt}\end{center}

\section{18. Thom's Influence on Algebraic
Geometry}\label{ch1958renethom:thoms-influence-on-algebraic-geometry}

\subsection{18.1 The Thom--Porteous
Formula}\label{ch1958renethom:the-thomporteous-formula-1}

As described in Section 12, the Thom--Porteous formula is essential in
enumerative geometry. It allows the computation of the number of
solutions to certain geometric problems using Chern class calculations.

\subsection{18.2 Intersection Theory}\label{ch1958renethom:intersection-theory}

Thom's work on the \textbf{homology of algebraic varieties} informed the
development of \textbf{intersection homology} (Goresky--MacPherson) and
the theory of characteristic classes of singular varieties.

\subsection{18.3 The Thom--Mather Theory and Algebraic
Geometry}\label{ch1958renethom:the-thommather-theory-and-algebraic-geometry}

Stratification theory is essential in algebraic geometry for studying
the topology of algebraic varieties with singularities. The
\textbf{Whitney stratification} of an algebraic variety and the
\textbf{Thom--Mather isotopy theorem} provide the foundation for the
topological study of singular algebraic varieties.

\begin{center}\rule{0.5\linewidth}{0.5pt}\end{center}

\section{19. Thom's Philosophy of
Mathematics}\label{ch1958renethom:thoms-philosophy-of-mathematics}

\subsection{19.1 Mathematical Realism}\label{ch1958renethom:mathematical-realism}

Thom was a philosophical \textbf{realist} about mathematics. He believed
that mathematical objects have an objective existence independent of
human thought. This position informed his approach to topology and
singularity theory.

\subsection{19.2 Qualititative
Mathematics}\label{ch1958renethom:qualititative-mathematics}

Thom advocated for a \textbf{qualitative} approach to mathematics,
emphasizing visualization and geometric intuition over algebraic
formalism. He saw topology as the study of ``forms'' and ``qualitative
structures.''

\subsection{19.3 The SMF and Mathematical
Politics}\label{ch1958renethom:the-smf-and-mathematical-politics}

Thom was active in the \textbf{Société Mathématique de France} and wrote
extensively on the philosophy and pedagogy of mathematics. He was
critical of the Bourbaki school's formalistic approach, arguing for a
return to geometric and intuitive methods.

\subsection{19.4 Semiophysics}\label{ch1958renethom:semiophysics}

In his later years, Thom developed \textbf{semiophysics}, an attempt to
apply his ideas about structural stability and morphogenesis to the
study of meaning, language, and cognition. His book \emph{Esquisse d'une
Sémiophysique} (1988) presents his vision.

\begin{center}\rule{0.5\linewidth}{0.5pt}\end{center}

\section{20. Thom's Work on the Classification of
Singularities}\label{ch1958renethom:thoms-work-on-the-classification-of-singularities}

\subsection{20.1 The Classification
Theorem}\label{ch1958renethom:the-classification-theorem-1}

Thom's classification of singularities of differentiable maps up to
\textbf{right-equivalence} is one of his important achievements.
For a map \(f: \mathbb{R}^n \to \mathbb{R}^p\), the classification
depends on: - The \textbf{codimension} of the singularity. - The
\textbf{determinacy} degree. - The \textbf{modality} (continuous
parameters in the classification).

\subsection{20.2 Simple, Unimodal, and Multimodal
Singularities}\label{ch1958renethom:simple-unimodal-and-multimodal-singularities}

Thom distinguished: - \textbf{Simple singularities:} Finitely many
classes with no continuous parameters (\(A_k\), \(D_k\), \(E_6\),
\(E_7\), \(E_8\)). - \textbf{Unimodal singularities:} Families depending
on one continuous parameter. - \textbf{Multimodal singularities:}
Families depending on two or more parameters.

\subsection{20.3 The Relation to Lie
Theory}\label{ch1958renethom:the-relation-to-lie-theory}

The classification of simple singularities (\(A_k\), \(D_k\),
\(E_{6,7,8}\)) corresponds exactly to: - The Dynkin diagrams of simple
Lie algebras. - The classification of \textbf{McKay quivers} for finite
subgroups of \(SL(2, \mathbb{C})\). - The classification of
\textbf{Kleinian singularities} \(\mathbb{C}^2/G\).

\begin{center}\rule{0.5\linewidth}{0.5pt}\end{center}

\section{21. Thom's Legacy in Modern
Mathematics}\label{ch1958renethom:thoms-legacy-in-modern-mathematics}

\subsection{21.1 Impact on Topology}\label{ch1958renethom:impact-on-topology}

Thom's work underpins much of modern topology: - \textbf{Cobordism}
(Milnor, Wall, Novikov). - \textbf{Bordism theories} (Brown--Peterson,
MU, MSO). - \textbf{Morse theory} (Floer homology, gauge theory). -
\textbf{Characteristic classes} (Chern--Weil theory, the Hirzebruch
signature theorem).

\subsection{21.2 Impact on Differential
Topology}\label{ch1958renethom:impact-on-differential-topology}

\begin{itemize}
\tightlist
\item
 \textbf{Transversality} and \textbf{Sard's theorem}.
\item
 \textbf{Intersection theory} and the \textbf{Alexander duality}.
\item
 \textbf{The h-principle} (Gromov, Eliashberg).
\end{itemize}

\subsection{21.3 Impact on Singularity
Theory}\label{ch1958renethom:impact-on-singularity-theory}

\begin{itemize}
\tightlist
\item
 \textbf{Arnold's classification} of singularities of functions and
 caustics.
\item
 \textbf{The theory of stable mappings} (Mather, Boardman).
\item
 \textbf{The theory of discriminants} (Givental, Vassiliev).
\end{itemize}

\subsection{21.4 Impact on Applied
Mathematics}\label{ch1958renethom:impact-on-applied-mathematics}

\begin{itemize}
\tightlist
\item
 \textbf{Catastrophe theory} (Zeeman, Poston, Stewart).
\item
 \textbf{Bifurcation theory} and \textbf{dynamical systems}.
\item
 \textbf{The theory of phase transitions} and \textbf{critical
 phenomena}.
\end{itemize}

\subsection{21.5 Impact on Algebraic
Geometry}\label{ch1958renethom:impact-on-algebraic-geometry}

\begin{itemize}
\tightlist
\item
 \textbf{Hironaka's resolution of singularities} uses stratification
 theory.
\item
 \textbf{Intersection homology} (Goresky--MacPherson) extends Poincaré
 duality to singular spaces.
\item
 \textbf{Motivic cohomology} uses Thom spaces in the construction of
 the \(\mathbb{A}^1\)-homotopy category.
\end{itemize}

\begin{center}\rule{0.5\linewidth}{0.5pt}\end{center}

\section{22. Thom's Philosophical
Writings}\label{ch1958renethom:thoms-philosophical-writings}

\subsection{22.1 The Éloge de
l'Intuition}\label{ch1958renethom:the-uxe9loge-de-lintuition}

Thom delivered a famous lecture ``Éloge de l'Intuition'' (In Praise of
Intuition) in which he argued for the primacy of geometric intuition
over formal algebraic reasoning in mathematical discovery. He saw the
Bourbaki program, despite its rigor, as potentially stifling to
mathematical creativity.

\subsection{22.2 Mathematics as a Science of
Forms}\label{ch1958renethom:mathematics-as-a-science-of-forms}

For Thom, mathematics was substantially the \textbf{science of forms}
(morphology). His view was that: - The advanced truths are
geometric. - Algebra is a language for describing geometric structures.
- Topology studies the most general qualitative properties of forms.

\subsection{22.3 The Crisis in Mathematical
Education}\label{ch1958renethom:the-crisis-in-mathematical-education}

Thom wrote critically about the ``modern mathematics'' reform movement
in education, which emphasized set theory and formal logic at the
expense of computation and geometry. He favored a return to intuitive,
geometric approaches to teaching.

\subsection{22.4 Semiophysics}\label{ch1958renethom:semiophysics-1}

Thom's later work on \textbf{semiophysics} attempted to apply
topological and dynamical ideas to linguistics and cognitive science. He
viewed meaning as arising from the \textbf{structural stability} of
mental representations, analogous to the stability of mathematical forms
under deformation.

\begin{center}\rule{0.5\linewidth}{0.5pt}\end{center}

\section{22. Awards and Recognition}\label{ch1958renethom:awards-and-recognition}

\textbf{1958:} - Fields Medal (International Mathematical Union) ---
awarded at the ICM in Edinburgh for cobordism theory, the transversality
theorem, and the Thom isomorphism.

\textbf{1970:} - Brouwer Medal (Dutch Mathematical Society).

\textbf{1976:} - John von Neumann Lecture Prize (Society for Industrial
and Applied Mathematics).

\textbf{1990:} - Honorary Member, London Mathematical Society.

\textbf{Elected to:} - French Academy of Sciences. - Royal Society of
London (Foreign Member). - American Academy of Arts and Sciences. -
Russian Academy of Sciences. - German Academy of Sciences Leopoldina.

\textbf{Doctor Honoris Causa:} - University of Oxford. - University of
Pisa. - University of Amsterdam. - University of Giessen.

\begin{center}\rule{0.5\linewidth}{0.5pt}\end{center}

\section{23. Selected Major
Publications}\label{ch1958renethom:selected-major-publications}

\begin{enumerate}
\def\labelenumi{\arabic{enumi}.}
\item
 \textbf{``Espaces fibrés en sphères et carrés de Steenrod.''}
 \emph{Annales Scientifiques de l'École Normale Supérieure} 69,
 109--182, 1952. --- \emph{Thom's thesis on spherical fibrations and
 Steenrod squares.}
\item
 \textbf{``Quelques propriétés globales des variétés
 différentiables.''} \emph{Commentarii Mathematici Helvetici} 28,
 17--86, 1954. --- \emph{Cobordism theory and the Thom transversality
 theorem.}
\item
 \textbf{``Les classes caractéristiques de Stiefel--Whitney.''}
 \emph{Séminaire Bourbaki} 5, 149--176, 1958. --- \emph{Characteristic
 classes and Steenrod operations.}
\item
 \textbf{``Les singularités des applications différentiables.''}
 \emph{Annales de l'Institut Fourier} 6, 43--87, 1956. ---
 \emph{Foundations of singularity theory.}
\item
 \textbf{``Stabilité structurelle et morphogenèse.''} \emph{W.A.
 Benjamin}, 1972. --- \emph{Catastrophe theory: classification of the
 seven elementary catastrophes.}
\item
 \textbf{``Topological methods in biology.''} \emph{Topology} 8,
 313--335, 1969. --- \emph{Applications of catastrophe theory to
 biology.}
\item
 \textbf{``Modèles mathématiques de la morphogenèse.''}
 \emph{Collection 10--18}, 1974. --- \emph{Mathematical models of
 morphogenesis.}
\item
 \textbf{``Esquisse d'une sémiophysique.''} \emph{InterÉditions}, 1988.
 --- \emph{Semiophysics: the application of topological ideas to
 language and meaning.}
\item
 \textbf{``Thom--Mather theory''} (with J. Mather). \emph{Stratified
 sets and the isotopy theorem.} Unpublished notes, 1960s.
\item
 \textbf{``Paraboles et catastrophes.''} \emph{Flammarion}, 1983. ---
 \emph{Collected essays on the philosophy of mathematics and science.}
\end{enumerate}

\begin{center}\rule{0.5\linewidth}{0.5pt}\end{center}

\begin{center}\rule{0.5\linewidth}{0.5pt}\end{center}

\section*{Exercises}
\addcontentsline{toc}{section}{Exercises}

\begin{exercise} Classify all smooth compact connected surfaces without boundary up to diffeomorphism. \end{exercise}

\begin{exercise} Show that two 3-manifolds are cobordant if and only if they are orientation-preserving diffeomorphic. \end{exercise}

\begin{exercise} Verify that the cusp catastrophe $V(x) = x^4 + ax^2 + bx$ has exactly two fold lines and one cusp point. \end{exercise}

\begin{exercise} Compute the cobordism ring $\\Omega_*^{SO}$ in low dimensions. \end{exercise}

\begin{center}\rule{0.5\linewidth}{0.5pt}\end{center}

\section*{Summary}
\addcontentsline{toc}{section}{Summary}

Thom created cobordism theory, showing that manifolds can be classified by algebraic invariants, and developed catastrophe theory as a mathematical framework for studying abrupt changes in continuous systems. His work connected topology, differential geometry, and applications to the natural sciences.

\section{24. References and Further
Reading}\label{ch1958renethom:references-and-further-reading}

\subsection{Official Sources}\label{ch1958renethom:official-sources}

\begin{itemize}
\tightlist
\item
 IMU Fields Medal 1958 Announcement:
 https://www.mathunion.org/imu-awards/fields-medal/fields-medals-1958
\item
 IHÉS: René Thom Biography.
\item
 Royal Society: Biographical Memoir.
\end{itemize}

\subsection{Key Technical References}\label{ch1958renethom:key-technical-references}

\begin{itemize}
\tightlist
\item
 Thom, R. \emph{Collected Works} (2 vols.). Springer, 1978.
\item
 Milnor, J. \emph{Characteristic Classes}. Princeton University Press,
 1974.
\item
 Arnold, V.I. \emph{Catastrophe Theory}. Springer, 1992.
\item
 Mather, J. ``Stability of \(C^\infty\) mappings.'' \emph{Annals of
 Mathematics} 87, 1968.
\end{itemize}

\subsection{Textbooks and Surveys}\label{ch1958renethom:textbooks-and-surveys}

\begin{itemize}
\tightlist
\item
 Bröcker, T. and tom Dieck, T. \emph{Kobordismen-Theorie}. Springer,
 1970.
\item
 Stong, R. \emph{Notes on Cobordism Theory}. Princeton University
 Press, 1968.
\item
 Hirsch, M. \emph{Differential Topology}. Springer, 1976.
\item
 Kosinski, A. \emph{Differential Manifolds}. Academic Press, 1993.
\end{itemize}

\subsection{Catastrophe Theory}\label{ch1958renethom:catastrophe-theory-1}

\begin{itemize}
\tightlist
\item
 Poston, T. and Stewart, I. \emph{Catastrophe Theory and Its
 Applications}. Pitman, 1978.
\item
 Zeeman, C. \emph{Catastrophe Theory: Selected Papers}. Addison-Wesley,
 1977.
\item
 Arnold, V.I. \emph{Theory of Catastrophes}. Springer, 1986.
\end{itemize}

\subsection{Singularity Theory}\label{ch1958renethom:singularity-theory-1}

\begin{itemize}
\tightlist
\item
 Arnold, V.I., Gusein-Zade, S.M., and Varchenko, A.N.
 \emph{Singularities of Differentiable Maps} (2 vols.). Birkhäuser,
 1985.
\item
 Dimca, A. \emph{Singularities and Topology of Hypersurfaces}.
 Springer, 1992.
\item
 Gibson, C.G. \emph{Singular Points of Smooth Mappings}. Pitman, 1979.
\end{itemize}

\subsection{Biography and Exposition}\label{ch1958renethom:biography-and-exposition}

\begin{itemize}
\tightlist
\item
 Hopf, H. ``The Work of René Thom.'' ICM 1958 Laudatio.
\item
 Sullivan, D. ``René Thom's work on geometric homology and bordism.''
 \emph{Bulletin of the AMS} 41, 2004.
\item
 Ghys, E. ``René Thom, catastrophe theorist, in the spotlight.''
 \emph{Maths News}, 2023.
\end{itemize}

\begin{center}\rule{0.5\linewidth}{0.5pt}\end{center}


\begin{center}\rule{0.5\linewidth}{0.5pt}\end{center}

\section*{See Also}
\addcontentsline{toc}{section}{See Also}

This chapter relates to the following work by other Fields Medalists:
\begin{itemize}
  \item \textbf{\hyperref[ch1986simondonaldson:simon-k.-donaldson-fields-medal-1986-a-comprehensive-report]{See Chapter on related work}} --- Thom's cobordism relates to Donaldson.
  \item \textbf{\hyperref[ch1966stephensmale:stephen-smale-fields-medal-1966-a-comprehensive-report]{See Chapter on related work}} --- Smale's work on h-cobordism relates to Thom.
\end{itemize}

\emph{(End of report)}
